\documentclass[a4paper,12pt]{article}
\usepackage[utf8]{inputenc}
\usepackage[spanish]{babel}
\usepackage{graphicx}
\usepackage{geometry}
\geometry{a4paper, margin=2.5cm}
\usepackage{hyperref}
\usepackage{titlesec}
\usepackage{enumitem}

\title{
    \vspace{-2cm}
    \begin{center}
        \textbf{\Large UNIVERSIDAD NACIONAL DE QUILMES} \\
        \vspace{1cm}
        % Si tienes el logo, puedes descomentar la siguiente línea
        % \includegraphics[width=0.4\textwidth]{logo.png} \\
        \vspace{1cm}
        \textbf{\Large TRABAJO DE INVESTIGACIÓN} \\
        \vspace{1.5cm}
        \textbf{Propuesta: Reconocimiento Robusto de la Lengua de Señas Argentina (LSA): Evolución de la Clasificación Estática a la Detección de Objetos en Tiempo Real} \\
    \end{center}
    \vspace{1.5cm}
}
\author{
    \begin{tabular}{c}
        \textbf{PRESENTADO POR:} \\
        JULIAN SANTIAGO GONZALEZ AVENDAÑO \\
        AGUSTIN GARCIA SMITH \\
        \vspace{1cm} \\
        \textbf{DIRIGIDO A:} \\
        INGENIERA ROXANA MARTINEZ \\
    \end{tabular}
}
\date{
    \vspace{2cm}
    \textbf{DEPARTAMENTO DE CIENCIA Y TECNOLOGIA} \\
    \vspace{0.5cm}
    \textbf{2026-1}
}

\begin{document}

\maketitle
\thispagestyle{empty}
\newpage

\section{Introducción y Antecedentes}

A diferencia de los enfoques tradicionales de aprendizaje automático que utilizan datos estructurados, este proyecto se enfoca en el campo de la Visión por Computadora empleando Aprendizaje Profundo. La meta principal es crear y entrenar una arquitectura de Inteligencia Artificial capaz de analizar imágenes y entender su significado semántico, centrándose en particular en la \textbf{Lengua de Señas Argentina (LSA)}.

El desarrollo de este modelo no solo tiene una finalidad académica, sino que aborda una problemática real de inclusión y accesibilidad para la comunidad sorda en Argentina. Establecer herramientas precisas de reconocimiento es el primer paso hacia sistemas de traducción automática.

\subsection{El Problema de Investigación}

Las investigaciones iniciales en este campo suelen abordar el reconocimiento de señas como un problema de \textit{clasificación de imágenes}, donde se asume que la imagen de entrada ya está perfectamente recortada y centrada sobre la mano. Sin embargo, en aplicaciones del mundo real (como una cámara web), el sistema recibe una imagen completa con diversos fondos, iluminación y ruido visual. 

En su trabajo sobre el reconocimiento de la LSA, Quiroga et al. (2017) evalúan diversas Redes Neuronales Convolucionales (CNN) sobre el dataset local LSA16 y concluyen que, si bien las CNNs son eficaces, \textbf{la pre-segmentación (recorte de las manos del fondo) es un factor crítico para aumentar significativamente la precisión del modelo}\footnote{Quiroga, F., et al. (2017). \textit{A Study of Convolutional Architectures for Handshape Recognition applied to Sign Language}. SEDICI.}.

Para complementar esta base, investigaciones recientes sobre el reconocimiento de lenguaje de señas y detección de objetos aplicada a gestos (ej., trabajos de los últimos cinco años) han empezado a adoptar arquitecturas de detección de objetos (\textit{one-stage detectors}) como YOLO, destacando su eficiencia en visión por computadora y aplicaciones de accesibilidad mediante IA.

Tomando esto como punto de partida, nos preguntamos: \textbf{¿De qué manera mejora la aplicabilidad y robustez del reconocimiento del Lenguaje de Señas al transicionar de un modelo de clasificación de imágenes estáticas (CNN) a un modelo avanzado de detección de objetos en tiempo real (YOLO) capaz de localizar espacialmente la mano en entornos no controlados?}

\subsection{Contribución Principal}
La contribución principal de este trabajo radica en la \textbf{evaluación de la robustez frente a fondos complejos} mediante el análisis comparativo experimental entre un modelo de clasificación tradicional (CNN) y uno de detección espacial (YOLO). Como aportes secundarios metodológicos, se destaca el desarrollo de un pipeline automático para la generación de anotaciones (bounding boxes) a partir de máscaras, y el análisis de factibilidad para aplicaciones de detección en tiempo real.

\section{Objetivos}

\subsection{Objetivo General}
Desarrollar y evaluar un sistema de reconocimiento automático para la Lengua de Señas Argentina (LSA) que demuestre la evolución y las ventajas de utilizar detección espacial robusta (YOLO) por sobre la clasificación estática clásica (CNN) en imágenes sin recortar.

\subsection{Objetivos Específicos}
\begin{itemize}
    \item \textbf{Preprocesamiento Avanzado:} Desarrollar un algoritmo basado en procesamiento clásico de imágenes (OpenCV) para generar automáticamente anotaciones espaciales (cajas delimitadoras o \textit{bounding boxes}) a partir de las máscaras de segmentación provistas en el dataset LSA16.
    \item \textbf{Línea de Base (Baseline):} Implementar y entrenar un modelo clasificador convolucional (CNN) utilizando el framework PyTorch, alimentado con imágenes pre-recortadas de las manos.
    \item \textbf{Detección de Objetos:} Implementar un modelo de la familia YOLO (\textit{You Only Look Once}) empleando \textit{Transfer Learning}, entrenado con las imágenes completas y las anotaciones espaciales autogeneradas.
    \item \textbf{Comparación y Análisis:} Contrastar cualitativa y cuantitativamente el desempeño de ambas arquitecturas para lidiar con el fondo de la imagen, discutiendo las implicancias metodológicas para el mundo real.
\end{itemize}

\section{Metodología}

La metodología propuesta se estructura en las siguientes fases:

\subsection{Dataset y Preprocesamiento}
El trabajo se basará en el dataset **LSA16** de la Universidad Nacional de La Plata, el cual contiene 800 imágenes de 16 clases de señas estáticas (alfabeto manual). Debido al volumen limitado de datos (Rios et al., 2020)\footnote{Rios, G., et al. (2020). \textit{Generación automática de videos utilizando Deep Learning. Aplicación en reconocimiento de gestos dinámicos}. SEDICI.}, se aplicarán rigurosas técnicas de \textit{Data Augmentation} (rotación, traslación, variación de brillo). 

Como el dataset provee imágenes \textit{Raw} y máscaras segmentadas (sin bounding boxes), se desarrollará un script metodológico para extraer las coordenadas de los cuadros delimitadores desde la máscara y exportarlas a formato YOLO.

\subsection{Arquitecturas a Comparar y Detalles Técnicos}
\begin{enumerate}
    \item \textbf{Modelo Clasificador (CNN):} Se diseñará una CNN convencional conformada por múltiples capas convolucionales y \textit{max pooling}, seguidas de capas densas (Fully Connected). Este modelo servirá como demostración de las vulnerabilidades al procesar imágenes no recortadas.
    \item \textbf{Modelo de Detección (YOLO):} Se adaptará una red preentrenada YOLOv8. Esta red no solo clasificará la seña, sino que predecirá las coordenadas $(x, y, ancho, alto)$ de la mano dentro del cuadro (regresión y clasificación en un solo paso).
\end{enumerate}

Para garantizar un enfoque experimental riguroso, se definirán los siguientes parámetros y consideraciones técnicas:
\begin{itemize}
    \item \textbf{Hiperparámetros y Entrenamiento:} Se empleará un optimizador estándar (ej. Adam o SGD) con un \textit{learning rate} inicial adecuado (ej. $1e-3$ o $1e-4$), utilizando un \textit{tamaño de batch} de 16 o 32 imágenes. Los modelos se entrenarán por una cantidad fija de \textit{épocas} (ej. 50-100), implementando \textit{Early Stopping} para mitigar el sobreajuste.
    \item \textbf{Estrategia de Validación:} El conjunto de datos será dividido sistemáticamente en subconjuntos de Entrenamiento, Validación y Prueba (ej. 70\%-15\%-15\%).
    \item \textbf{Hardware Utilizado:} Los entrenamientos se ejecutarán sobre entornos con aceleración por GPU (e.g., Google Colab con NVIDIA T4/A100 o equivalente local) para asegurar tiempos de procesamiento viables.
\end{itemize}

\subsection{Criterios de Evaluación}
Para el clasificador se utilizarán métricas estándar: \textit{Accuracy, Precision, Recall, F1-Score} y \textit{Matriz de Confusión}. Para YOLO, el desempeño se medirá a través de la \textit{Intersection over Union (IoU)} y \textit{mean Average Precision (mAP)}.

\section{Resultados Experimentales Esperados}

Se espera demostrar empíricamente que el modelo de clasificación disminuye fuertemente su precisión al introducir fondos complejos. Por el contrario, se espera que YOLO logre aislar la característica de interés y mantener un alto \textit{mAP}.

Para respaldar el análisis comparativo y de métricas, el documento final presentará:
\begin{itemize}
    \item \textbf{Curvas de Entrenamiento:} Gráficas de pérdida (\textit{Loss}) y métricas a lo largo de las épocas.
    \item \textbf{Tablas Comparativas:} Rendimiento global de ambos modelos, contrastando precisión y tiempo de inferencia.
    \item \textbf{Matrices de Confusión:} Para visualizar el desempeño por clase y detectar confusiones comunes entre gestos.
    \item \textbf{Ejemplos Visuales de Detección:} Comparación visual de imágenes con cajas delimitadoras detectadas frente al \textit{ground truth}.
    \item \textbf{Análisis de Errores:} Discusión detallada sobre los casos difíciles, como fondos con colores similares a la piel u oclusiones.
\end{itemize}

Adicionalmente, el proyecto dejará documentado un flujo de trabajo replicable para convertir máscaras de segmentación en anotaciones de objetos espaciales.

\section{Conclusiones Preliminares y Trabajo Futuro}

Esta investigación sentará las bases para comprender la necesidad de integrar la localización espacial en el reconocimiento de gestos. Como trabajo futuro, se proyecta escalar el sistema desde imágenes estáticas (LSA16) hacia secuencias de video con gestos dinámicos, utilizando el dataset de mayor complejidad \textbf{LSA64}, requiriendo para ello arquitecturas espaciotemporales (como Redes Neuronales Recurrentes o CNNs 3D).

\end{document}
